\chapter{Hearing Tests}

\section*{Definition and Overview}
Hearing tests are procedures used to assess how well a person hears and to identify the type, degree, and cause of hearing loss. They range from simple screening checks, such as newborn hearing screening and whispered voice tests, to comprehensive audiological assessments performed by audiologists. Hearing tests measure how loud sounds need to be before they are heard, how clearly speech is understood, and how well the outer, middle, and inner parts of the ear function. They provide the information needed to diagnose hearing loss and to fit and program hearing aids effectively.

\section*{Epidemiology}
Hearing loss affects around one in six Australians, and testing is therefore a common part of health care. Universal newborn hearing screening tests more than 98\% of babies born in Australia, identifying congenital hearing loss in about one to three per 1,000 births. Adults are increasingly tested as hearing loss prevalence rises with age, and workplace noise-exposure programs require periodic audiometric testing for workers in high-noise industries. Despite this, many people with significant hearing loss have never had a hearing test, and delayed diagnosis remains common.

\section*{Aetiology and Pathophysiology}
\begin{itemize}
    \item \textbf{Purpose:} hearing tests characterise hearing function across the outer ear, middle ear, inner ear (cochlea), and auditory nerve
    \item \textbf{Pure-tone audiometry:} measures the quietest sounds (thresholds) heard at different frequencies, producing an audiogram
    \item \textbf{Speech testing:} assesses the ability to hear and understand words, reflecting real-world listening
    \item \textbf{Immittance testing:} tympanometry and acoustic reflexes evaluate middle ear function and the reflex pathway
    \item \textbf{Otoacoustic emissions:} measure the response of outer hair cells of the cochlea, used mainly in newborns and young children
    \item \textbf{Auditory brainstem response:} records electrical activity from the auditory nerve and brainstem, useful when behavioural testing is not possible
\end{itemize}

\section*{Clinical Features}
\begin{itemize}
    \item Tests are indicated when a person struggles in conversation, turns up the volume, or is told they do not hear well
    \item Children may be tested for speech delay, inattention, or recurrent ear infections
    \item Newborns undergo routine screening before leaving hospital
    \item Workers in noisy industries undergo surveillance audiometry
    \item Tinnitus, dizziness, or a family history of hearing loss may prompt testing
    \item Sudden hearing loss requires urgent testing and medical review
\end{itemize}

\section*{Differential Diagnosis}
\begin{itemize}
    \item Test results must be interpreted in context: a single screening test is not the same as a full audiological assessment
    \item Temporary threshold shifts after noise exposure can mimic permanent loss and resolve with rest
    \item Wax, fluid, or infection in the middle ear can produce conductive patterns that may resolve with treatment
    \item Children who are difficult to test may give unreliable results, requiring repeat or alternative testing
    \item Functional (non-organic) hearing loss occurs when test responses do not match observed behaviour
\end{itemize}

\section*{When to Seek Medical Care}
Hearing tests are appropriate for anyone who suspects hearing loss, for children with speech or listening concerns, and as part of routine screening for newborns and older adults. A doctor or audiologist can advise which tests are needed.

\textbf{Call 000 (or local emergency services) for:}
\begin{itemize}
    \item Sudden hearing loss, especially with dizziness, ringing, or fullness in one ear
    \item Hearing loss with severe headache, facial weakness, or other neurological symptoms
    \item Hearing loss following head injury
\end{itemize}

\section*{Diagnosis and Investigations}
\begin{itemize}
    \item \textbf{Otoscopy:} examination of the ear canal and eardrum before formal testing
    \item \textbf{Pure-tone audiometry:} performed in a sound booth using headphones or insert earphones, mapping thresholds from low to high frequencies
    \item \textbf{Speech audiometry:} measures the level at which speech is detected and understood
    \item \textbf{Tympanometry:} assesses eardrum movement and middle ear pressure
    \item \textbf{Otoacoustic emissions and auditory brainstem response:} objective tests for infants and complex cases
    \item \textbf{Newborn screening:} automated otoacoustic emissions or auditory brainstem response testing in the first weeks of life
    \item \textbf{Imaging:} CT or MRI when structural causes such as otosclerosis or acoustic neuroma are suspected
\end{itemize}

\section*{Management}
\begin{itemize}
    \item \textbf{Interpretation:} an audiologist or audiometrist explains the results and their implications
    \item \textbf{Referral:} results guide referral to an ear, nose, and throat specialist where medical or surgical causes are found
    \item \textbf{Hearing aid fitting:} the audiogram is used to program amplification for the individual's pattern of loss
    \item \textbf{Follow-up testing:} hearing is re-tested periodically to monitor change and adjust devices
    \item \textbf{Surveillance audiometry:} workers in noisy environments are tested at regular intervals to detect early damage
    \item \textbf{Education:} results inform advice on hearing protection, communication strategies, and rehabilitation
\end{itemize}

\section*{Complications}
\begin{itemize}
    \item Untested and untreated hearing loss leads to isolation, depression, and cognitive decline
    \item Normal results can provide false reassurance if a person's difficulties are situational or fluctuating
    \item Incorrectly fitted aids based on inadequate testing give poor outcomes
    \item Missed congenital hearing loss delays language development and education
    \item Hearing loss progression may go undetected without regular follow-up testing
\end{itemize}

\section*{Prognosis and Outcomes}
Early, accurate testing allows timely intervention, and outcomes are consistently better when hearing loss is identified and managed promptly. Newborn screening followed by early intervention enables children with congenital hearing loss to develop language skills at near-normal rates. For adults, testing that leads to effective amplification improves communication, social participation, mood, and quality of life. Regular testing also monitors progressive conditions so that management can be adjusted before the loss becomes disabling.

\section*{Prevention and Self-Care}
\begin{itemize}
    \item Have hearing checked as part of routine health care, especially from middle age
    \item Ensure newborns receive their routine hearing screen
    \item Seek testing promptly if friends or family comment on hearing difficulties
    \item Participate in workplace audiometric surveillance if exposed to noise
    \item Re-test at recommended intervals to monitor any change
    \item Protect hearing between tests by using hearing protection in loud environments
    \item Report sudden changes in hearing immediately
\end{itemize}

\section*{Sources and Further Reading}
The following reputable sources provide further information for healthcare students and patients seeking reliable, current advice about hearing tests.

\begin{itemize}
    \item Healthdirect. \textit{Hearing tests and what they involve.} https://www.healthdirect.gov.au/hearing-tests
    \item Hearing Australia. \textit{About hearing tests.} https://www.hearing.com.au/
    \item Australian Government Department of Health. \textit{Newborn hearing screening program.} https://www.health.gov.au/
    \item NHS. \textit{Tests for hearing loss.} https://www.nhs.uk/conditions/hearing-loss/
    \item Mayo Clinic. \textit{Hearing tests: what to expect.} https://www.mayoclinic.org/
    \item MedlinePlus. \textit{Hearing tests.} https://medlineplus.gov/ency/article/003341.htm
\end{itemize}
